\section{IA para la ciencia: descubrimiento científico guiado por LLM}

\subsection{De la regresión simbólica a la evolución guiada por LLM}

\begin{knowledgebox}{Problema de regresión simbólica}
La \textbf{regresión simbólica} (Symbolic Regression) es la tarea de descubrir ecuaciones a partir de datos observacionales empíricos. Por ejemplo, dados datos de medición del movimiento de Marte, descubrir la tercera ley de Kepler.

Los métodos tradicionales (como la programación genética PySR) mantienen una población de expresiones candidatas, que evoluciona mediante mutación y cruce. Los LLM pueden potenciar este proceso aportando conocimiento científico previo.
\end{knowledgebox}

\subsection{LaSR: regresión simbólica guiada por conceptos}

\begin{importantbox}{Marco LaSR (Language-guided Symbolic Regression)}
LaSR infiere conjuntamente hipótesis de programa y conceptos de alto nivel:

$$P(\pi, C \mid D) \propto P(D \mid \pi) \cdot P(\pi \mid C) \cdot P(C)$$

donde:
\begin{itemize}
    \item $P(D \mid \pi)$: la verosimilitud del programa $\pi$ respecto a los datos $D$ (calculada ejecutando el programa)
    \item $P(\pi \mid C)$: la probabilidad del programa dado el concepto $C$ (el conocimiento previo del LLM)
    \item $P(C)$: la distribución previa de los conceptos (el conocimiento científico del LLM)
\end{itemize}
\end{importantbox}

Flujo de trabajo de LaSR:
\begin{enumerate}
    \item Mantener una población de programas (ecuaciones) candidatos
    \item \textbf{Evolución de hipótesis}: mutación simbólica clásica + mutación guiada por LLM (basada en la biblioteca de conceptos)
    \item \textbf{Evaluación de aptitud}: qué tan bien se ajusta el programa a los datos
    \item \textbf{Abstracción de conceptos}: el LLM abstrae el conjunto de programas en descripciones de alto nivel en lenguaje natural (como «crecimiento/decaimiento exponencial», «ley de potencia»)
    \item \textbf{Evolución de conceptos}: el LLM combina conceptos existentes para generar conceptos nuevos (como «crecimiento exponencial» + «dependencia de la temperatura» $\to$ «distribución de Boltzmann»)
    \item La biblioteca de conceptos retroalimenta la evolución de hipótesis, formando un ciclo cerrado
\end{enumerate}

\subsection{Resultados experimentales}

\begin{itemize}
    \item En la tarea de descubrimiento de las ecuaciones de Feynman, LaSR supera tanto a la programación genética pura como a la versión sin aprendizaje de conceptos
    \item Incluso usando modelos de lenguaje locales (no de frontera), logra superar a la programación genética pura
    \item Los prompts del usuario (hints) pueden aportar una ganancia adicional
    \item La expresión de la ley de Coulomb que descubrió LaSR solo necesita 4 pasos de simplificación para llegar a la forma estándar, mientras que la expresión equivalente obtenida por PySR requiere una simplificación considerable
\end{itemize}

\subsection{Nuevo descubrimiento: leyes de escalamiento de LLM}

Para verificar que el método no se limita a «memorizar fórmulas conocidas», el equipo usó LaSR para descubrir una nueva ley de escalamiento de LLM que incorpora el parámetro del \textbf{número de ejemplos few-shot}:

\begin{warningbox}{Hallazgo interesante}
Un número grande de ejemplos few-shot resulta perjudicial para los modelos de baja capacidad, pero una vez que la capacidad del modelo supera cierto umbral, más shots producen un mejor desempeño. Al combinar este hallazgo con la ley de Chinchilla se obtuvo una ley de escalamiento más precisa.
\end{warningbox}

\subsection{Resumen de la sección}

\begin{itemize}
    \item La evolución guiada por LLM es una herramienta poderosa para la ciencia empírica
    \item La capacidad de abstracción del LLM (descubrimiento de conceptos) puede potenciar aún más el proceso de descubrimiento
    \item Se puede extender a entradas visuales de alta dimensión (usando VLM)
    \item Retos abiertos: validación de hipótesis, representación de conceptos más allá del lenguaje natural, y escalamiento a espacios de búsqueda mayores
\end{itemize}
